\documentclass[blue]{beamer}

\mode<presentation> {
  \setbeamertemplate{background canvas}[vertical shading][bottom=red!10,top=blue!10]

  \usetheme{Boadilla}
  \usefonttheme[onlysmall]{serif}
}


\usepackage{pgf,pgfarrows,pgfnodes,pgfautomata,pgfheaps,pgfshade}
\usepackage{amsmath,amssymb,bm}
\usepackage{colortbl}
\usepackage[english]{babel}
\usepackage[latin1]{inputenc}


\setbeamercovered{dynamic}


\def\R{\mathbb{R}}
\def\Q{\mathbb{Q}}
\def\N{\mathbb{N}}
\def\E{\mathbb{E}}
\newcommand{\ps}[2]{ \langle #1 , #2 \rangle }
\def\1{\mathbf{1}}
\newcommand{\nor}[1]{\left\| #1 \right\|}
\def\leq{\leqslant}
\def\geq{\geqslant}
\def\et{\quad \textrm{and} \quad}
\def\ou{\quad \textrm{or} \quad}
\def\boite{\hskip 0.5mm {\Box} \hskip 0.5mm}


\DeclareMathOperator\PPE{PPE}
\DeclareMathOperator\vNM{vNM}
\DeclareMathOperator\argmax{argmax}
\DeclareMathOperator\Prob{Prob}
\DeclareMathOperator\supp{supp}
\DeclareMathOperator\Eq{Eq}
\DeclareMathOperator\inte{int}
\DeclareMathOperator\co{co}
\DeclareMathOperator\cl{cl}
\DeclareMathOperator\Dirac{Dirac}
\DeclareMathOperator\F{F}
\DeclareMathOperator\As{As}
\DeclareMathOperator\Ir{Ir}
\DeclareMathOperator\MRS{MRS}
\DeclareMathOperator\Range{Im}
\DeclareMathOperator\DIV{DIV}

%------------------------------ theoren styles

\theoremstyle{definition}
\newtheorem*{assumption}{Assumption}
\newtheorem*{conjecture}{Conjecture}

\theoremstyle{example}
\newtheorem*{proposition}{Proposition}
\newtheorem*{remark}{Remark}


%-----------------------------------------

\title[Corporate Finance]{Theory of Corporate Finance}
\author[Victor Filipe]{V. Filipe Martins-da-Rocha}
\institute[FGV EESP]{Escola de Economia de S\~ao Paulo}
\date[April--June, 2026]{Part 1.a. Stock Market Equilibrium: Objective of the Firm}
\subject{Economic Theory}


\begin{document}

\frame{\titlepage}

%\part<presentation>{Main Talk}

%------------------ outline

\begin{frame}\frametitle{Outline}
\begin{itemize}
\item Definition of a stock market equilibrium: firms and investors
\item Separation of ownership and control
    \begin{itemize}
    \item Fisher Separation Theorem
    \end{itemize}
\item The financial structure of the firm
    \begin{itemize}
    \item Irrelevance results: the Modigliani--Miller's theorems
    \item Taxation
    \item Limited liability and bankruptcy
    \end{itemize}
\end{itemize}
\end{frame}

%------------------------------------------------

\begin{frame}\frametitle{A simple intertemporal model}

\begin{block}{Time}
Two dates: $t\in \{0,1\}$
\end{block}

\begin{block}{Uncertainty at $t=1$}
Finite set $S$ of states of nature
\end{block}

\begin{block}{Commodity}
One divisible commodity available for trade at every period and state
\end{block}

\end{frame}

%------------------------------------------------

\begin{frame}\frametitle{A simple intertemporal model}

\begin{block}{Investors}
A finite set $I$ of investors
\end{block}

\begin{block}{Firms}
One (for simplicity) firm privately owned by investors
\end{block}

\begin{block}{Financial assets}
A finite set $A$ of financial contracts (promises)
\end{block}

\end{frame}


%------------------------------------------------

\begin{frame}\frametitle{The firm}

\begin{block}{Production set}
The production possibilities of the firm is represented by a set
\[
Y \subset \R \times \R^S
\]
\end{block}

\begin{block}{Production plan}
A production plan for the firm is a vector in $Y$
\[
y = (y_0,y_1) \quad \textrm{where} \quad y_1=(y_s)_{s\in S}
\]
\end{block}


\end{frame}

%------------------------------------------------

\begin{frame}\frametitle{Interpretation}

\begin{assumption}
    \begin{itemize}
    \item instantaneous production is impossible
    \item the firm invests at $t=0$
    \item the firm produces at $t=1$
    \end{itemize}
In other words,
\[
Y \subset \R_{-} \times \R^S_{+}
\]
\end{assumption}

Consider a vector
\[
y=(y_0,(y_s)_{s\in S}) \in Y \subset \R_{-} \times \R^S_{+}
\]

\begin{block}{At $t=0$}
$y_0 \leq 0$ means that $-y_0$ units of the good are used as input at $t=0$
\end{block}

\begin{block}{At $t=1$ and state $s$}
$y_s \geq 0$ means that $y_s$ units of good $\ell$ are produced at state $s$
\end{block}


\end{frame}

%------------------------------------------------

\begin{frame}\frametitle{Firms are privately owned}

\begin{itemize}
\item We denote by $\xi^i \geq 0$ the share of the firm owned by investor $i$ at $t=0$
\item Shares are percentages of ownership, in particular we must have
\[
\sum_{i\in I} \xi^i =1
\]
\item Some investor $i$ might have no shareholding in the firm, i.e., $\xi^i =0$
\item Some investor $i$ might be the single owner of the firm, i.e., $\xi^i=1$
\end{itemize}

\end{frame}

%------------------------------------------------

\begin{frame}\frametitle{Competitive stock markets}

At $t=0$, there is a competitive market for stocks of the firm
    \begin{itemize}
    \item the price is linear
    \item investors take prices as given
    \item the price of the firm in units of the unique commodity is denoted by $E$ (equity)
    \end{itemize}
If agent $i$ wants to buy the share $\alpha \geq 0$, then he should pay $E \alpha$ units of the unique commodity


\end{frame}

%------------------------------------------------

\begin{frame}\frametitle{Financial assets or securities}

A financial asset $a \in A$ is a claim to a state contingent bundle
\[
\kappa(a) =(\kappa_s(a))_{s\in S} \in \R^S
\]

\begin{block}{Portfolios}
Each agent can choose a portfolio $\theta \in \R^A$ at $t=0$
\begin{itemize}
\item if $\theta(a) \geq 0$, the agent expects to receive at $t=1$ and contingent to state $s$ the amount $\kappa_s(a) \theta(a)$ of the commodity
\item if $\theta(a) \leq 0$, the agent promises to deliver at $t=1$ and contingent to state $s$ the amount $\kappa_s(a) \theta(a)$ of the commodity
\end{itemize}
\end{block}

\begin{block}{Riskless (or risk-free, or non-contingent) bond}
An asset $a$ is called risk-free if the vector $\kappa(a)$ is not random, e.g.,
\[
\forall s\in S, \quad \kappa_s(a) = 1
\]
\end{block}

\end{frame}

%------------------------------------------------

\begin{frame}\frametitle{Competitive security markets}

At $t=0$, there is a competitive market for securities
    \begin{itemize}
    \item the price is linear
    \item investors take prices as given
    \item investor can buy ($\theta(a)>0$) or sell ($\theta(a)<0$) any security $a\in A$
    \item the price of security $a$ in units of the unique commodity is denoted by $q(a)$
    \end{itemize}
If agent $i$ wants to buy the portfolio $\theta \in \R^A$, then he should pay (or receive)
\[
q \cdot \theta = \sum_{a \in A} q(a) \theta(a)
\]
 units of the unique commodity

\end{frame}

%------------------------------------------------

\begin{frame}\frametitle{Competitive security markets: interest rates}

\begin{itemize}
\item Denote by $w(a)$ the resources (in units of the unique commodity) spent by some agent in asset~$a$
\item It is equivalent to buy $\theta(a)$ units of the asset where
\[
w(a) = q(a) \theta(a)
\]
\item The payoff at $t=1$ is given by
\[
\kappa_s(a) \theta(a)
\]
\item Equivalently, the (stochastic) gross return on asset~$a$, defined by
\[
\kappa_s(a) \theta(a) = R_s(a) w(a)
\]
implying that the net return $r_s(a)$ defined by $R_s(a) = 1 + r_s(a)$ satisfies
\[
r_s(a) = \frac{\kappa_s(a) - q(a)}{q(a)}
\]
\end{itemize}
\end{frame}


%------------------------------------------------

\begin{frame}\frametitle{Investors}

\begin{block}{Characteristics}
\begin{itemize}
\item an initial endowment $\omega^i_0 \geq 0$ of the commodity
\item an initial vector of shares $\xi^i \in \R_+$
\item a state contingent endowment $\omega^i_1=(\omega^i_s)_{s\in S} \in \R^S_+$
\item a Bernoulli function $u^i : \R_+ \to [-\infty,\infty)$
\item a discount factor $\beta^i >0$
\item beliefs about uncertainty represented by $\nu^i \in \Prob(S)$
\end{itemize}
\end{block}

\end{frame}


%------------------------------------------------

\begin{frame}\frametitle{Investors}

\begin{block}{Actions}
At $t=0$, agent $i$ chooses a consumption $c_0 \in \R_+$ and makes plans for future consumption $c_1\in \R^S_+$
\end{block}

\begin{block}{Expected utility}
\[
U^i(c_0,c_1) = u^i(c_0) + \beta^i \sum_{s\in S} \nu^i_s u^i(c_s)
\]
\end{block}

\end{frame}

%------------------------------------------------

\begin{frame}\frametitle{Firms' actions}

At $t=0$, given security price $q$, the firm
    \begin{itemize}
    \item chooses a (physical) production plan $$y=(y_0,y_1) \in Y$$
    \item chooses a financial plan $\beta \in \R^A$
        \begin{itemize}
        \item $\beta(a) >0 $ means the firm is selling $\beta(a)$ units of security $a$
        \item $\beta(a) < 0 $ means the firm is purchasing $\beta(a)$ units of security $a$
        \end{itemize}
    \end{itemize}
This production-finance plan $(y,\beta)$ leads to a time and state dependent cash flow
\[
\delta = (\delta_0,\delta_1) \in \R \times \R^S
\]
where
\[
\delta_0 = y_0 + q \cdot \beta \et \delta_s = y_s - \kappa_s\cdot \beta
\]

\end{frame}


%------------------------------------------------

\begin{frame}\frametitle{Firms' actions at $t=0$}

Owners of the firm at $t=0$ (agents $i$ with $\xi^i >0$) must finance the expenditure $\delta_0$:
    \begin{itemize}
    \item if $\delta_0 <0$, agent $i$ has to deliver $\xi^i \delta_0$ units of the commodity
    \item if $\delta_0 > 0$, agent $i$ receives $\xi^i \delta_0$ units of the commodity
    \end{itemize}

\end{frame}


%------------------------------------------------

\begin{frame}\frametitle{Firms' actions at $t=1$ contingent to state $s$}

Owners of the firm at $t=1$ (agents $i$ with $\alpha^i >0$) receive (or deliver) the surplus $\delta_1$:
    \begin{itemize}
    \item if $\delta_s <0$, agent $i$ has to deliver $\alpha^i \delta_s$ units of the commodity in state $s$
    \item if $\delta_s > 0$, agent $i$ receives $\alpha^i \delta_s$ units of the commodity in state $s$
    \end{itemize}

\end{frame}

%------------------------------------------------

\begin{frame}\frametitle{Usual interpretation at $t=0$}

\begin{itemize}
\item Firm $j$ chooses a production plan $$y=(y_0,y_1) \in Y$$ where $y_0 \leq 0$ is an investment and $y_s \geq 0$ is a production
\item Firm $j$ decides to finance its investment $y_0$ by short-selling assets, i.e., choosing $\beta \in \R^A_+$
    \begin{itemize}
    \item One may have
    \[
    y_0 + q \cdot \beta \leq 0
    \]
    in that case only a part of the investment is financed by the firm and the rest should be finance by shareholders
    \item One may have
    \[
    y_0 + q \cdot \beta \geq 0
    \]
    in that case all the investment is financed by the firm and shareholders may even receive a surplus
    \end{itemize}
\end{itemize}

\end{frame}


%------------------------------------------------

\begin{frame}\frametitle{Usual interpretation at $t=1$}

When the state $s$ realizes at $t=1$
\begin{itemize}
\item the firm produces $y_s$ units of the commodity
\item the firm must honor its financial commitment by paying its debt $\kappa_s \cdot \beta$
\item each agent $i$ receives (or delivers) the residual $\delta_s$
    \begin{itemize}
    \item if $y_s > \kappa_s \cdot \beta$ then agent $i$ receives $\delta_s$ units of the commodity
    \item if $y_s <\kappa_s \cdot \beta$ then agent $i$ delivers $-\delta_s$ units of the commodity
    \end{itemize}
\end{itemize}

\end{frame}

%------------------------------------------------

\begin{frame}\frametitle{Types of ownership structure}

\begin{block}{Sole proprietorship}
\begin{itemize}
\item There is no stock market for the firm
\item The firm is individually owned:
\[
\exists i \in I, \quad \xi^{i}=1
\]
\item The owner manages the firm in his own interests
\end{itemize}
\end{block}

\begin{block}{Partnership}
\begin{itemize}
\item There is no stock market for the firm
\item Several shareholders
\item Possible disagreement among partners on the objective to be assigned to the firm
\end{itemize}
\end{block}

\end{frame}

%------------------------------------------------

\begin{frame}\frametitle{Types of ownership structure}

\begin{block}{Corporation}
\begin{itemize}
\item There is a stock market
\item Separation of ownership and control
\end{itemize}
\end{block}

\end{frame}


%------------------------------------------------

\begin{frame}\frametitle{Limited liability}

Consider that the firm can only finance its production plan through the riskless asset $a_0$ (non-contingent bond) characterized by the claim
\[
\forall s\in S, \quad \kappa_s(a_0) = 1
\]
\begin{itemize}
\item At $t=0$, the firm decides to issue debt by selling the portfolio $Z \1_{\{a_0\}}$

\item In that case, the firm is said \emph{levered}

\item Shareholders may be protected and have \emph{limited liability} in the sense that they will never deliver

\item In other words, if $\delta_s<0$, no shareholder will pay the debt to creditors

\item In that case a bankruptcy rule applies and the whole production plan $y_s$ is distributed among creditors
\end{itemize}
\end{frame}

%------------------------------------------------

\begin{frame}\frametitle{Investors' budget sets}

Each investor $i$ takes as given (and perfectly anticipates):
    \begin{itemize}
    \item the firm's production-finance plan $(y,\beta)$
    \item the security price vector $q$ and the equity price~$E$
    \end{itemize}
and chooses
    \begin{itemize}
    \item a consumption plan $(c_0,c_1) \in \R_+ \times \R^S_+$
    \item a security portfolio $\theta\in \R^A$ and a shareholdings vector $\alpha \in \R_+$
    \end{itemize}
such that ...

\end{frame}

%------------------------------------------------

\begin{frame}\frametitle{Investors' budget sets}

\begin{itemize}
\item the budget restriction at $t=0$
\begin{equation}\label{eq:budget-0}
c_0 + q\cdot \theta + E \alpha \leq \omega^i_0 + (E + \delta_0) \xi^i
\end{equation}
\item the budget restriction at each $s\in S$,
\begin{equation}\label{eq:budget-s}
c_s \leq \omega^i_s + \kappa_s \cdot \theta + \delta_s \alpha
\end{equation}
\end{itemize}

The set of actions $(c,\theta,\alpha)$ satisfying (\ref{eq:budget-0}) and (\ref{eq:budget-s}) is denoted by $B^i(q,E)$
\begin{remark}
Observe that we should write $\delta_0(\alert{q,y,\beta})$ and $B^i(q,E,\alert{y,\beta})$
\end{remark}

\end{frame}

%------------------------------------------------

\begin{frame}\frametitle{Stock market Equilibrium}

A stock market equilibrium is a list of
\[
\left\{ (y,\beta), (q,E), (c^i,\theta^i,\alpha^i)_{i\in I}\right\}
\]
of
\begin{itemize}
\item production-finance plan $(y,\beta)$ for the firm
\item security and equity prices $(q,E)$
\item action $(c^i,\theta^i,\alpha^i)$ of each investor $i$
\end{itemize}
such that
\end{frame}

%------------------------------------------------

\begin{frame}\frametitle{Stock market Equilibrium}

\begin{itemize}
\item[(a)] investor $i$'s action is optimal, i.e.,
\[
(c^i,\theta^i,\alpha^i) \in \argmax \{ U^i(c) \ \colon \ (c,\theta,\alpha) \in B^i(q, E) \}
\]
\item[(b)] security and stock markets clear, i.e.,
\[
\forall a\in A, \quad \beta(a) = \sum_{i\in I} \theta^i(a) \et \sum_{i\in I} \alpha^i =1
\]
\item[(c)] commodity markets clear, i.e.,
\[
\sum_{i\in I} c^i = \sum_{i\in I} \omega^i + y
\]
\end{itemize}

\end{frame}


%------------------------------------------------

\begin{frame}\frametitle{The objective of the firm}

\alert{
When choosing the production-finance plan, what is the firm's objective?
}

\end{frame}

%------------------------------------------------

\begin{frame}\frametitle{The objective of the firm: an arbitrary solution}

A possible \alert{arbitrary} answer is: to maximize the firm's value $V$ defined by
\begin{eqnarray*}
V & = & E + \delta_0 \\
 & = & E + q \cdot \beta + y_0 \\
 &  = & E + D + y_0
\end{eqnarray*}
which depends on
    \begin{itemize}
    \item the firm's decision $(y,\beta)$
    \item the associated equilibrium prices $(E,q)$
    \end{itemize}

\end{frame}

%------------------------------------------------

\begin{frame}\frametitle{The objective of the firm: Separation of ownership and control}

\begin{itemize}
\item The firm belongs to investors having a positive share
\item Agent $i$ with $\xi^i>0$ would like the manager of the firm to choose the production-finance plan $(y^i,\beta^i)$ that maximizes his expected utility
\item There is a problem if shareholders disagree about which production-finance plan the firm should choose
\item If it happens that there exists a function (objective) such that all agents agree that the firm's choice should be to maximize this function, then we have unanimity among shareholders
\item If there is unanimity, we can separate \textcolor{red}{control} from \textcolor{red}{ownership} and employ a manager to maximize the objective
\end{itemize}

\end{frame}

%------------------------------------------------

\begin{frame}\frametitle{Separation of ownership and control}

\begin{itemize}
\item Under which condition can we separate \textcolor{red}{control} from \textcolor{red}{ownership}?
    \begin{itemize}
    \item if there is no uncertainty
    \item if markets are complete
    \end{itemize}
\item Actually, under the two previous conditions, it happens that for each agent $i$, the optimal production-finance plan $(y^i,\beta^i)$ according to agent~$i$'s preferences coincides with the
plan maximizing the (present) value $V$ of the firm
\item In that case, shareholders can delegate control of the firm to a manager whose goal is to maximize the firm's value
\end{itemize}

\end{frame}

%------------------------------------------------

\begin{frame}\frametitle{Literature}

\begin{thebibliography}{8}

\begin{small}

  \beamertemplatearticlebibitems

    \bibitem{Fisher30}
    Fisher, I.
    \newblock {\em The Theory of Interest}
    \newblock New York: A. M. Kelley, 1930

    \bibitem{Dreze74}
    Dr\`eze, J. H.
    \newblock Investment under private ownership: optimality, equilibrium and stability.
    \newblock {\em Allocation Under Uncertainty: Equilibrium and Optimality}, Mew York: Wiley, 1974

    \bibitem{GrossmanHart79}
   Grossman, S. J. and Hart, O. D.
    \newblock A theory of competitive equilibrium in stock market economies.
    \newblock {\em Econometrica} 47, 1979

    \bibitem{GrossmanStiglitz80}
   Grossman, S. J. and Stiglitz, J.
    \newblock Stockholder unanimity in making production and financial decisions.
    \newblock {\em Quarterly J. of Economics} 94, 1980

\end{small}


    \end{thebibliography}


\end{frame}

%------------------------------------------------

\begin{frame}\frametitle{Recent developments}


\begin{small}

\begin{thebibliography}{8}


  \beamertemplatearticlebibitems


    \bibitem{Geanaloplosetali}
    Geanakoplos, J., Magill, M., Quinzii, M. and Dr\`eze, J.
    \newblock Generic inefficiency of stock market equilibrium when markets are incomplete.
    \newblock {\em Journal of Mathematical Economics} 19, 1990

   \bibitem{KelseyMilne}
   Kelsey, D. and Milne, F.
    \newblock The existence of equilibrium in incomplete markets and the objective function of the firm.
    \newblock {\em Journal of Mathematical Economics} 25, 1996

    \bibitem{TvedeCres}
   Tvede, M. and Cr\'es, H.
    \newblock Voting in assemblies of shareholders and incomplete markets.
    \newblock {\em Economic Theory} 26, 2005

    \bibitem{MagillQuinzii08JME}
    Magill, M. and Quinzii, M.
    \newblock Normative properties of stock market equilibrium with moral hazard
    \newblock {\em Journal of Mathematical Economics} 44, 2008


    \end{thebibliography}

\end{small}
\end{frame}

%------------------------------------------------

\begin{frame}\frametitle{Recent developments}

\begin{small}

\begin{thebibliography}{8}


  \beamertemplatearticlebibitems

      \bibitem{MagillQuinzii09}
    Magill, M. and Quinzii, M.
    \newblock The probability approach to general equilibrium with production
    \newblock {\em Economic Theory} 39, 2009


     \bibitem{Herings09}
    Herings, J.~J.
    \newblock {\em The theory of the firm: bargaining and competitive equilibrium}
    \newblock  SAET Conference, Ischia, 2009

\bibitem{CarcelespovedaCoenpirani09}
    Carceles-Poveda, E. and Coen-Pirani, D.
    \newblock Shareholder's unanimity with incomplete markets
    \newblock {\em International Economic Review} 50, 2009



    \end{thebibliography}

\end{small}

\end{frame}


%------------------------------------------------

\begin{frame}\frametitle{Recent developments}

\begin{thebibliography}{8}


  \beamertemplatearticlebibitems



     \bibitem{MagillQuinzii08}
    Magill, M. and Quinzii, M.
    \newblock {\em A co-moment criterion for the choice of risky investment by firms}
    \newblock  International Economic Review, 2010

    \bibitem{Dierker09}
    Dierker, E. and Dierker, H.
    \newblock {\em Dr\`eze equilibria and welfare maxima}
    \newblock Economic Theory, 2011

    \bibitem{BejanBidian10}
    Bejan, C. and Bidian, F.
    \newblock Ownership structure and efficiency in large economies
    \newblock {\em Economic Theory}, 2011

    \end{thebibliography}

\end{frame}

%------------------------------------------------

\begin{frame}\frametitle{Standard assumptions}


For each investor $i$,
\begin{itemize}
\item[(C.1)] the function $u^i$ is continuously differentiable, concave and strictly increasing
\item[(C.2)] endowments in the commodity are strictly positive, i.e., $\omega^i_0>0$ and $\omega^i_s >0$ for each $s$
\item[(C.3)] every state $s$ has a positive probability, i.e., $\nu^i_s >0$ for each $s\in S$
\item[(C.4)] individually rational consumption plans are interior, i.e.,
\[
U^i(c_0,c_1) \geq U^i(\omega^i_0, \omega^i_1) \Longrightarrow c_0 >0 \et c_s >0, \quad \forall s\in S
\]
\end{itemize}

\end{frame}

%------------------------------------------------

\begin{frame}\frametitle{Fundamental asset pricing theorem}


\begin{proposition}
Consider that Assumptions C.1 to C.3 are satisfied and assume that
\[
\left\{ (y,\beta), (q,E), (c^i,\theta^i,\alpha^i)_{i\in I}\right\}
\]
is a stock market equilibrium.
There exists a vector $$\lambda=(\lambda_s)_{s\in S} \in \R^S_{++}$$ of state price deflators such that
\[
\forall a \in A, \quad q(a) = \sum_{s\in S} \lambda_s \kappa_s(a)
\]
and
\[
E = \sum_{s\in S} \lambda_s \delta_s = \sum_{s\in S} \lambda_s [ y_s - \kappa_s \cdot \beta]
\]
\end{proposition}

\end{frame}


%------------------------------------------------

\begin{frame}\frametitle{Personalized state prices}


\begin{proposition}
Consider that Assumptions C.1 to C.4 are satisfied and assume that
\[
\left\{ (y,\beta), (q,E), (c^i,\theta^i,\alpha^i)_{i\in I}\right\}
\]
is a stock market equilibrium.
We have
\[
\forall a \in A, \quad q(a) = \sum_{s\in S} \MRS^i_s \kappa_s(a)
\]
and for every shareholder $i$ of the firm at $t=1$ (i.e., $\alpha^i>0$), we have
\[
E = \sum_{s\in S} \MRS^i_s \delta_s \quad \textrm{where} \quad \MRS^i_s\alert{(c^i)} := \beta^i \nu^i_s \frac{ \partial u^i(c^i_s)}{\partial u^i(c^i_0)}
\]
\end{proposition}
If agents have homogenous expectations, i.e., $\nu^i = \nu$, then $\beta^i \frac{ \partial u^i(c^i_s)}{\partial u^i(c^i_0)}$ is called a stochastic discount factor (with respect to $\nu$)

\end{frame}

%------------------------------------------------

\begin{frame}\frametitle{Impact of the firm's decision on investors}

Fix a stock market equilibrium
\[
\left\{ (y,\beta), (q,E), (c^i,\theta^i,\alpha^i)_{i\in I}\right\}
\]
and consider the following ``out-of-equilibrium'' path argument:
\begin{itemize}
\item Change the firm's decision by a small amount $(\Delta y, \Delta \beta)$
\item Neglect the impact of this change on equilibrium prices other than the firm's stock price
\item Consumer $i$'s evaluation of the impact of this change in his expected utility will depend on agent $i$'s presumed effect
\[
\Delta^i E
\]
on the firm's equity price
\end{itemize}

\end{frame}

%------------------------------------------------

\begin{frame}\frametitle{Impact of the firm's decision on investor $i$}

We denote by $\Delta^i U^i$ the change on his utility presumed by agent $i$ as a consequence of the firm's change $(\Delta y, \Delta \beta)$
\begin{eqnarray*}
\frac{\Delta^i U^i}{\partial u^i(c^i_0)} & = & \xi^i \left[ \Delta^i E + \Delta \delta_0\right]  \\
& + & \alpha^i \left[ \sum_{s\in S} \MRS^i_s \Delta \delta_s - \Delta^i E \right]
\end{eqnarray*}
where
\[
\Delta \delta_0 = \Delta y_0 + q \cdot \Delta \beta
\]
and
\[
\Delta \delta_s = \Delta y_s - \kappa_s \cdot \Delta \beta
\]

\end{frame}

%------------------------------------------------

\begin{frame}\frametitle{Competitive price perceptions}

Grossman and Hart (1980) propose to assume that investors have \alert{{\em competitive price perceptions}}

\begin{block}{Competitive price perceptions}
Each agent $i$ evaluates the value at $t=0$ of any income stream at $t=1$ using his own marginal rates of substitution as state price deflators
\end{block}
In other words, if $w$ is a vector in $\R^S$, then agent $i$ evaluates the value at $t=0$ of the claim to this income stream as follows
\begin{equation}\label{eq:competitive-price-perceptions}
\sum_{s\in S} \MRS^i_s w_s
\end{equation}
\begin{itemize}
\item Assume that an asset paying the stream $w$ is introduced in the economy after the equilibrium has been reached
\item At equilibrium (i.e., given his consumption plan $c^i$), agent~$i$ is indifferent between buying or selling this new asset if the price is given by~(\ref{eq:competitive-price-perceptions})
\end{itemize}

\end{frame}

%------------------------------------------------

\begin{frame}\frametitle{Market consistency}

If the income stream $w$ is marketed, i.e.,
\[
\exists \theta_w \in \R^A, \quad w_s = \kappa_s \cdot \theta_w, \ \forall s\in S
\]
then agent $i$'s evaluation is consistent with the market price of $w$, i.e.,
\[
\sum_{s\in S} \MRS^i_s w_s =  q \cdot \theta_w
\]
In particular, even if two agents have different marginal rates of substitution, they will agree on the present value of the future income stream $w$


\end{frame}

%------------------------------------------------

\begin{frame}\frametitle{Impact of the firm's decision on investor $i$}

Assuming agent $i$ has competitive price perceptions, then
\[
\Delta^i E = \sum_{s\in S} \MRS^i_s \Delta \delta_s
\]
and
\[
\Delta^i U^i = \partial u^i(c^i_0) \xi^i \left[ \sum_{s\in S} \MRS^i_s \Delta \delta_s + \Delta \delta_0\right]
\]
\begin{block}{Present value of the firm's dividend stream}
With competitive price perceptions, a stockholder would welcome any production-finance policy of the firm which maximizes the \alert{\textit{present value of the firm's dividend stream evaluated at his
own marginal rate of substitution} }
\end{block}

\end{frame}

%------------------------------------------------

\begin{frame}\frametitle{Disagreement}

If  two agents $i_1\neq i_2$ have different marginal rates of substitution, i.e.,
\[
(\MRS^{i_1}_s)_{s\in S} \neq (\MRS^{i_2}_s)_{s\in S}
\]
then a change $(\Delta y, \Delta \beta)$ of the firm's policy may lead to
\[
\Delta^{i_1} U^{i_1} >0 \et \Delta^{i_2} U^{i_2} <0
\]
In that case shareholders do not agree on a common objective
\end{frame}

%------------------------------------------------

\begin{frame}\frametitle{Unanimity}

If security markets are complete, i.e.,
\[
\Range \kappa = \R^S
\]
then there exists a unique state price deflator $\lambda \in \R^S_{++}$ and all agents have the same competitive price perception, i.e.,
\[
\forall i\in I, \quad \MRS^i_s = \lambda_s
\]
In that case, there is unanimity among agents about the firm's objective: \alert{maximize its value}
\begin{eqnarray*}
V & = & \delta_0 + \sum_{s\in S} \lambda_s \delta_s \\
& = & y_0 + \sum_{s\in S} \lambda_s y_s \\
& = & E + D + y_0
\end{eqnarray*}

\end{frame}

%------------------------------------------------

\begin{frame}\frametitle{Unanimity II}

Observe that if for any feasible production change $\Delta y$, there exist $\Delta \theta \in \R^A$ and $\Delta \alpha \in \R$ such that
\[
\Delta y_s = \sum_{a\in A} \kappa_s(a) \Delta \theta(a) + y_s \Delta \alpha
\]
then there is unanimity among agents about the firm's objective

\end{frame}

%------------------------------------------------

\begin{frame}\frametitle{Strong unanimity}

We denote by $D$ the set of possible no-default production-finance plans $(y,\beta)$ available to the firm, i.e.,
\[
D \equiv \left\{ (y,\beta) \in Y \times \R^A \ \colon \ \forall s\in S, \quad y_s \geq \kappa_s \cdot \beta \right\}
\]
Given a family $\lambda=(\lambda_s)_{s\in S} \in \R^S_{++}$ of stochastic discount factors, we let $B^i(q,\lambda,y,\beta)$ be the budget set associated to the security price $q$ and the equity price
\[
\widetilde{E}(\lambda,y,\beta)\equiv \sum_{s\in S} \lambda_s \{ y_s - \kappa_s \cdot \beta \}
\]

\end{frame}

%------------------------------------------------

\begin{frame}\frametitle{Strong unanimity}

\begin{theorem}[Strong unanimity]
Fix a competitive stock market equilibrium
\[
\{(k,\beta),(q,E),(c^i,\theta^i,\alpha^i)_{i\in I}\}
\]
where \alert{agents' marginal rates of substitution coincide}, i.e.,
\[
\exists \lambda=(\lambda_s)_{s\in S}, \quad \forall i\in I, \quad \MRS^i_s = \lambda_s.
\]
Assume that the firm has chosen $(y,\beta)$ to maximize its value defined by
\[
\forall (y^\prime,\beta^\prime) \in D, \quad V(y^\prime,\beta^\prime) \equiv y_0^\prime + q\cdot \beta^\prime + \sum_{s\in S} \lambda_s [y_s^\prime - \kappa_s \beta^\prime]
\]
For each investor~$i$, we have
\[
(y,\beta) \in \argmax \{ U^i(c') \ \colon \ (c',\theta',\alpha') \in B^i(q,\lambda,y',\beta') \ \textrm{and} \ (y',\beta') \in A \}
\]
\end{theorem}

\end{frame}

\end{document} 